% ==========================================================
% STANDARDIZED PREAMBLE (STATIC CORE) — DO NOT EDIT THIS PART
% ==========================================================

\documentclass[11pt,a4paper]{article}

\usepackage{iftex}
\ifPDFTeX
  \usepackage[utf8]{inputenc}
  \usepackage[T1]{fontenc}
  \usepackage{stix2}
\else
  \usepackage{fontspec}
  \usepackage{unicode-math}
  \setmainfont{STIX Two Text}
  \setmathfont{STIX Two Math}
\fi

\usepackage[margin=2.5cm]{geometry}
\usepackage{parskip}
\usepackage{microtype}
\usepackage{enumitem}

\usepackage{amsmath, amssymb}

\usepackage{tocloft}
\setlength{\cftbeforesecskip}{2pt}
\setlength{\cftbeforesubsecskip}{1pt}

\widowpenalty=10000
\clubpenalty=10000

\usepackage{xcolor}
\definecolor{myblue}{RGB}{0, 51, 153}

\usepackage{sectsty}
\partfont{\color{myblue}}
\chapterfont{\color{myblue}}
\sectionfont{\color{myblue}}
\subsectionfont{\color{myblue}}
\subsubsectionfont{\color{myblue}}

\usepackage{hyperref}

\renewcommand{\cftsecfont}{\bfseries\color{myblue}}
\renewcommand{\cftsubsecfont}{\color{myblue}}
\renewcommand{\cftsubsubsecfont}{\color{myblue}}
\renewcommand{\cftsecpagefont}{\bfseries}
\renewcommand{\cftsubsecpagefont}{}
\renewcommand{\cftsubsubsecpagefont}{}

\hypersetup{
  colorlinks=true,
  linkcolor=myblue,
  citecolor=myblue,
  urlcolor=myblue
}

\usepackage{titlesec}
\titlespacing{\section}{0pt}{12pt}{6pt}
\titlespacing{\subsection}{0pt}{10pt}{5pt}
\titlespacing*{\section}{0pt}{2.5ex plus 1ex minus .2ex}{1.5ex plus .2ex}
\titlespacing*{\subsection}{0pt}{2ex plus 0.8ex minus .2ex}{1.2ex plus .2ex}
\titlespacing*{\subsubsection}{0pt}{1.8ex plus 0.6ex minus .2ex}{1ex plus .2ex}

\makeatletter
\newcommand{\StartBlueDoc}[3][]{%
  \title{\vspace{-2em}\textbf{#2}%
    \if\relax\detokenize{#3}\relax\else\\[2pt]{\large #3}\fi
    \vspace{+0.6em}}
  \author{#1}
  \date{\vspace{-0.6em}\today}
  \begin{document}
  \maketitle
  \tableofcontents
  \vspace{1em}\hrule\vspace{1em}
}
\makeatother

% ==========================================================
% CUSTOMIZATION ZONE
% ==========================================================


\newcommand{\DocAuthor}{}
\newcommand{\DocTitle}{Testing Research Software}
\newcommand{\DocSubtitle}{}

\setcounter{tocdepth}{2}

\usepackage{booktabs}
\usepackage{tabularx}
\newcolumntype{L}{>{\raggedright\arraybackslash}X}
\newcolumntype{T}{>{\raggedright\arraybackslash}p{0.36\textwidth}}

% Header rows of the table "At a glance": a section, with its number and its
% title.
\newcommand{\GlanceSection}[1]{%
  \multicolumn{3}{@{}l}{\textbf{\ref{#1}\quad\nameref{#1}}}}

% Cite a principle of the other document by its title. A \ref cannot
% reach into another PDF, so the title is the reference.
\newcommand{\DesignRef}[1]{\emph{Design Principles}, ``#1''}

% ==================
% Start the document
% ==================
\StartBlueDoc[\DocAuthor]{\DocTitle}{\DocSubtitle}



% ==========================================================
\section*{About this guide}
\phantomsection
\addcontentsline{toc}{section}{About this guide}
% ==========================================================

This guide is about testing research software: showing that the code
computes what it claims to compute. Together with \emph{Design Principles
for Research Software}, cited below as \emph{Design Principles}, it forms one
set of guidelines for all code written for a research purpose, in academia or
in industry, in any programming language and scientific field. It is meant
for looking things up while working, not for reading from cover to cover.

A test runs code on an input and compares the result with an expected value.
The comparison is easy; knowing the expected value is hard. In research
software, the typical failure is not a crash but a plausible, wrong number,
and a test catches it only if its expected value is right and comes from
somewhere other than the code under test. A test is therefore only as good as
its oracle, the source of its expected value.
Section~\ref{sec:sources} is about where oracles come from,
Section~\ref{sec:writing} about writing tests, and
Section~\ref{sec:honest} about keeping a test suite honest as the code
changes.

Tests show that the thing is built right: that the code computes what the
method defines. Whether it is the right thing, that is, whether the method
fits the scientific question, is validation, a scientific question that no
test can answer (\DesignRef{Build the right thing before building the thing
right}).

Two related topics belong to \emph{Design Principles}. That code must be
tested, and which code, is stated in its principles ``Testing is not
optional'' and ``Test the whole workflow on a small dataset''. Checks that
the code makes on every run, such as input validation and assertions of
invariants, are not tests; they are the subject of ``Garbage in, useful error
out'' and ``Wear belt and braces''.



\subsection*{Normative language and the ladder}
\phantomsection
\addcontentsline{toc}{subsection}{Normative language and the ladder}
\label{sec:normative}

The words \emph{must}, \emph{should}, and \emph{may}, and the plain
imperatives ``prefer'' and ``avoid'', carry the meaning defined in the
subsection ``Normative language'' of \emph{Design Principles}. The rules of
this guide apply along the same ``Ladder of Research Software'': the
\emph{must}s are part of the floor from the second rung on, and what rises
with the rung is how many kinds of oracle the tests use, how strong they
are, and how much of the input space they reach. When rules conflict, the
priorities in \DesignRef{Get your priorities right} decide. Each rule is
stated in one of the two documents only.



\subsection*{Terminology}
\phantomsection
\addcontentsline{toc}{subsection}{Terminology}
\label{sec:terminology}

The terms defined in \emph{Design Principles}, such as package, library code,
workflow code, entry point, host stack, and caller, carry the same meaning
here. The terms of testing, such as test, test suite, oracle, and tolerance,
and the types of tests, such as unit test, end-to-end test, and snapshot
test, are defined in \emph{Types of Research Software Tests}, which also
describes what each type catches and what it misses. This guide uses them
with that meaning.



% The checklist gets a page of its own, so that it can be printed.
\clearpage
% ==========================================================
\section*{At a glance}
\phantomsection
\addcontentsline{toc}{section}{At a glance}
% ==========================================================

Every principle of this guide in one line. The number and the title link to
the principle.

{\small
\begin{tabularx}{\textwidth}{@{}lTL@{}}
  \toprule
  \GlanceSection{sec:sources} \\
  \midrule
  \ref{sec:oracle} & \nameref{sec:oracle}
    & Never take an expected value from the code under test. \\
  \ref{sec:known-answers} & \nameref{sec:known-answers}
    & Use exact results from analytic solutions, hand calculations, and
      certified references, and record their source. \\
  \ref{sec:invariants} & \nameref{sec:invariants}
    & Check properties that every correct result has, such as conservation
      laws, bounds, and normalizations. \\
  \ref{sec:metamorphic} & \nameref{sec:metamorphic}
    & Check how the output must change when the input changes in a known
      way. \\
  \ref{sec:recovery} & \nameref{sec:recovery}
    & Simulate data with known parameters, and check that the method
      recovers them with honest uncertainty. \\
  \ref{sec:convergence-order} & \nameref{sec:convergence-order}
    & Check that the error of a discretization falls at the rate of the
      method. \\
  \ref{sec:second-opinion} & \nameref{sec:second-opinion}
    & Compare with an independent implementation, ideally a naive one
      written for the test. \\
  \addlinespace
  \GlanceSection{sec:writing} \\
  \midrule
  \ref{sec:float-tests} & \nameref{sec:float-tests}
    & Compare within a tolerance derived from the numerical error of the
      method, not bit for bit. \\
  \ref{sec:flaky-tests} & \nameref{sec:flaky-tests}
    & Fix the random generator, and assert statistical claims with bounds
      that hold for any seed. \\
  \ref{sec:edges} & \nameref{sec:edges}
    & Test empty, single, tied, missing, and extreme inputs, and the
      documented behavior there. \\
  \ref{sec:generated-inputs} & \nameref{sec:generated-inputs}
    & Check invariants and relations on generated inputs, the same on every
      run of the suite. \\
  \ref{sec:small-tests} & \nameref{sec:small-tests}
    & Write mostly fast unit tests, a few end-to-end tests, and run slow
      tests separately. \\
  \ref{sec:isolated} & \nameref{sec:isolated}
    & No dependence on other tests, global state, or the network; no
      confidential data. \\
  \ref{sec:snapshots} & \nameref{sec:snapshots}
    & Store only validated outputs, compare them within a tolerance, and
      explain every update. \\
  \addlinespace
  \GlanceSection{sec:honest} \\
  \midrule
  \ref{sec:bug-tests} & \nameref{sec:bug-tests}
    & Reproduce each bug in a failing test before fixing it. \\
  \ref{sec:can-fail} & \nameref{sec:can-fail}
    & See every new test fail, for example by breaking the code on
      purpose. \\
  \ref{sec:skips} & \nameref{sec:skips}
    & Report every skipped test with its reason, and never silence a failing
      one. \\
  \ref{sec:test-code} & \nameref{sec:test-code}
    & Keep tests simple, name each for its claim, and make failures
      explain themselves. \\
  \bottomrule
\end{tabularx}}



% ==========================================================
\section{Where correctness comes from}
\label{sec:sources}
% ==========================================================

\subsection{A test is only as good as its oracle}
\label{sec:oracle}

A test compares a result with an expected value, and it is only as good as
the source of that value, its oracle. A test meant to show that a result is
correct must not take its expected value from the code under test. Running
the code and copying its output into the test makes the test pass by
construction; it is then a snapshot test (Subsection~\ref{sec:snapshots}),
which shows that the result has not changed since, not that it is right.

An oracle can be wrong too, and a test that agrees with a wrong oracle is
worse than no test, because it lends the wrong result authority. Prefer
oracles that are simpler than the code under test and can be checked by
reading them, and state where each comes from
(Subsection~\ref{sec:known-answers}).

No single kind of oracle catches every error. A known answer checks a few
points of the input space exactly; an invariant checks every input, but only
one property of the result; parameter recovery checks the whole method on
realistic data, but only within sampling error. Each entry point should be
tested against at least one oracle that is independent of its code, and the
more depends on a method, the more kinds of oracle should test it.
Table~\ref{tab:oracles} lists strong oracles for common kinds of code.

\begin{table}[h]
\centering
\small
\caption{Strong oracles for common kinds of code. The numbers refer to the
subsections of this guide.}
\label{tab:oracles}
\begin{tabularx}{\textwidth}{@{}LLL@{}}
  \toprule
  Code & Typical silent error & Strong oracles \\
  \midrule
  Numerical routines, such as transforms, special functions, and solvers of
  linear systems
    & A wrong result in part of the input range, such as for large
      magnitudes or nearly singular input
    & Known answers (\ref{sec:known-answers}), invariants
      (\ref{sec:invariants}), a naive implementation
      (\ref{sec:second-opinion}) \\
  \addlinespace
  Discretizations, such as solvers of differential equations and quadrature
  rules
    & Convergence at a lower order than the method has
    & The order of convergence and manufactured solutions
      (\ref{sec:convergence-order}), conservation laws
      (\ref{sec:invariants}) \\
  \addlinespace
  Statistical estimators and model fits
    & Biased estimates, or standard errors and intervals that are too narrow
    & Parameter recovery with coverage (\ref{sec:recovery}), certified
      reference datasets (\ref{sec:known-answers}), an independent
      implementation (\ref{sec:second-opinion}) \\
  \addlinespace
  Stochastic simulations
    & The right mean with the wrong distribution
    & Statistical bounds on distributions (\ref{sec:flaky-tests}), limiting
      cases with known answers (\ref{sec:known-answers}), invariants
      (\ref{sec:invariants}) \\
  \addlinespace
  Data transformations, such as parsing, joining, normalizing, and filtering
    & Rows dropped, duplicated, or misaligned with their identifiers
    & Tiny cases checked by hand (\ref{sec:known-answers}), reordering and
      duplication (\ref{sec:metamorphic}), counts (\ref{sec:invariants}) \\
  \addlinespace
  Workflows
    & Steps that do not fit together, or options that never reach their step
    & An end-to-end run on simulated input with a known result
      (\ref{sec:recovery}), counts and totals (\ref{sec:invariants}) \\
  \bottomrule
\end{tabularx}
\end{table}



\subsection{Know the answer before you ask}
\label{sec:known-answers}

The most direct oracle is a known answer: the exact result for a particular
input, obtained without the code under test. Sources include:

\begin{itemize}
  \item Analytic solutions and limiting cases, such as a model whose closed
    form is known for one value of a parameter, or the behavior of an
    estimator as the sample size grows.
  \item Tiny cases worked out by hand or with a computer algebra system, such
    as the likelihood of three observations or the alignment of two short
    sequences.
  \item Certified reference values, such as the Statistical Reference
    Datasets of NIST, whose results are certified to many digits.
  \item Published values from the literature or from another established
    tool.
\end{itemize}

State next to each expected value where it comes from: the derivation, the
hand calculation, or the reference, with the table and the number of digits
given. An expected value without a stated source cannot be checked by a
reviewer; it is a magic number (\DesignRef{No magic numbers}).

Make sure that the reference computes the same quantity. Published values
often rest on other conventions or defaults: a variance divided by $n$ or by
$n-1$, a rate or a scale parameterization, another rule for ties, or another
convergence tolerance (\DesignRef{Numerics and defaults are part of the
API}). A disagreement may then be a difference of conventions rather than a
bug. Derive the tolerance from the precision of the reference: a value
published with four significant digits cannot test the fifth.

Known answers are exact but sparse. They exist mostly for small or simple
inputs and so test the easy corner of the input space; the other kinds of
oracle in this section reach the rest.



\subsection{Test invariants, not only examples}
\label{sec:invariants}

Passing on a few happy-path inputs is a weak scientific test. Prefer
invariants: properties that every correct result has, whatever the input.
Examples are conservation laws, such as of mass, energy, or the total of the
counts; normalizations, such as probabilities that sum to one; bounds, such
as a nonnegative variance or a p-value between zero and one; symmetries,
such as a symmetric distance matrix; and consistency relations, such as row
and column totals that agree.

An invariant needs no known answer, so it can be checked on any input,
including large, random, and extreme ones where no known answer exists
(Subsection~\ref{sec:generated-inputs}). It checks only one property of the
result, however, and a result can have it and still be wrong: probabilities
that sum to one can still be the wrong probabilities. Combine invariants
with other oracles (Subsection~\ref{sec:oracle}).

Examples are useful, but they do not prove that the method is correct.
Encode the scientific claims the software makes and test those claims
directly. The invariants worth testing are often the ones worth asserting in
the code itself (\DesignRef{Wear belt and braces}): the test checks them on
the inputs someone chose, the assertion on the data that users bring.



\subsection{Change the input, predict the output}
\label{sec:metamorphic}

Where the correct result for an input is unknown, its relation to the result
for a related input often is known. Change the input in a way whose effect
on the output the method determines, and test that the output changes
accordingly. Such metamorphic relations include:

\begin{itemize}
  \item \textbf{Reordering.} Permuting the observations leaves an estimate
    unchanged, and permuting the rows of the input permutes the rows of the
    output in the same way, with the identifiers still aligned.
  \item \textbf{Rescaling.} Expressing the input in other units changes the
    result by the corresponding factor, and shifting the coordinates shifts
    an estimated position.
  \item \textbf{Duplicating and splitting.} Duplicating every observation
    leaves a maximum likelihood estimate unchanged and divides its standard
    error by $\sqrt{2}$. Splitting the data, computing an additive
    statistic such as a count or a sum on each part, and combining the parts
    gives the result for the whole.
  \item \textbf{Monotonicity.} A higher confidence level never gives a
    narrower interval, and a higher detection threshold never detects more
    events.
\end{itemize}

Such relations catch errors that leave each single result plausible, such as
two axes mixed up, identifiers misaligned after a sort, or a unit converted
twice. Like invariants, they need no known answer and can be checked on any
input (Subsection~\ref{sec:generated-inputs}). Derive each relation from the
method, not from the code. A relation that holds only approximately, for
example because reordering changes the rounding of a sum or because an
iterative method stops at a tolerance, is checked within a tolerance of its
own (Subsection~\ref{sec:float-tests}).



\subsection{Plant the answer, then find it}
\label{sec:recovery}

Statistical methods, and inverse problems in general, should be tested by
parameter recovery: simulate data from the model with known parameters,
apply the method, and check that it recovers them. For such code it is often
the strongest test, because it exercises the whole method on data of
realistic size and structure, where no hand-computed answer exists.

Check the uncertainty as well as the estimate. An estimator can find the
right value and still report standard errors or intervals that are too
narrow, and this is the more dangerous error, because it makes wrong
conclusions look certain. Over many simulated datasets, the estimates should
scatter around the true values as much as the reported standard errors say,
and the intervals should contain the true values at their nominal rate,
their coverage. Simulate under the null hypothesis as well: a test at the 5\%
level should then reject in about 5\% of the datasets. Include hard cases,
such as small samples, parameters near the boundary of their range, and weak
signals, because that is where estimators break first.

Assert recovery within bounds derived from the sampling error, with a known
and very small false-failure rate (Subsection~\ref{sec:flaky-tests}). The
simulator is the oracle, so it should be written from the definition of the
model, independently of the code of the method: a simulator and an estimator
that share a mistake, such as the same wrong parameterization, agree
perfectly. Studies of coverage need many replicates and are slow; run them
with the slow tests (Subsection~\ref{sec:small-tests}).



\subsection{Check the order of convergence}
\label{sec:convergence-order}

A discretization, such as a solver of differential equations, a quadrature
rule, or a finite-difference scheme, should be tested by its observed order
of convergence, not only by its error at one resolution. Run it at a
sequence of resolutions against an exact solution, estimate the order from
how fast the error falls, and compare it with the order the method has:
halving the step of a second-order method divides the error by about four.
Choose resolutions fine enough that the leading term of the error dominates
and coarse enough that rounding error does not.

The error at one resolution is a weak test. A bug that reduces a
second-order method to first order, such as a boundary condition applied one
step late, still gives a small error on a fine enough grid, and a tolerance
loose enough to pass at one resolution passes such bugs too. The order is a
sharp oracle: a single number known from the theory of the method, which
most bugs of a discretization change.

Where no exact solution exists, manufacture one. Choose a smooth function as
the solution, insert it into the equations to obtain the source terms and
boundary conditions it requires, add these to the problem, and test that the
solver converges to the chosen function at the expected order. The same idea
applies beyond grids: the error of a Monte Carlo estimate falls as $N^{-1/2}$
in the number of samples, and an iterative solver converges at the rate its
theory gives.



\subsection{Get a second opinion}
\label{sec:second-opinion}

Compare results with an independent implementation of the same method. The
best is often a naive one, written for the test: a plain loop instead of
vectorized code, a dense matrix instead of a sparse one, a direct sum
instead of a fast transform, or all permutations instead of a random sample
of them. It is slow, but simple enough to be checked by reading
(\DesignRef{KISS: Keep it simple, stupid}), and on small inputs it provides
an expected value for every input the fast code accepts. Keep it in the test
code, not in the package.

Another established package, possibly in another language, gives a second
opinion as well. Either call it in the test, with its version pinned in the
test environment, or store its results as known answers together with its
version (Subsection~\ref{sec:known-answers}); a new release of it can change
its numbers (\DesignRef{Numerics and defaults are part of the API}). Compare
the same quantity under the same conventions.

Agreement is evidence only as far as the implementations are independent.
Two implementations written from the same derivation, from the same paper,
or by the same person share their misreadings and agree on the wrong answer.
A disagreement does not say which implementation is wrong; resolve it with
another oracle before changing either.



% ==========================================================
\section{Writing tests}
\label{sec:writing}
% ==========================================================

\subsection{Compare floats with tolerance}
\label{sec:float-tests}

Tests should compare floating-point results within a tolerance, not bit for
bit. Bit-identical reference files treat platform noise as a logic bug and
make the test suite fragile.

Assert agreement within the tolerance that defines reproducibility
(\DesignRef{Once is luck, twice is science}). Use relative tolerances for
quantities of varying magnitude and absolute tolerances for quantities near
zero. Tolerances are constants in the sense of \DesignRef{No magic numbers}:
derive each from the expected numerical error of the method, rather than
adjusting it until the tests pass on one machine. A tolerance that is too
loose is as harmful as one that is too tight, because it passes wrong
results; it should be smaller than the smallest error that would matter
(Subsection~\ref{sec:can-fail}).



\subsection{Flaky tests cry wolf}
\label{sec:flaky-tests}

A test that fails at random teaches people to ignore failures, and a test
suite whose failures are ignored checks nothing. Tests of stochastic methods
must therefore fix the random generator (\DesignRef{Once is luck, twice is
science}), so that every run gives the same outcome.

A fixed generator makes a test repeatable, not convincing: a test that passes
for only one seed can hide a bug. Where the claim is statistical, such as an
unbiased estimate or a coverage rate, assert it with bounds whose
false-failure rate is known and very small, for example a deviation of
several standard errors, so that the test would pass for other seeds as well.
The bound is a constant in the sense of \DesignRef{No magic numbers}. Never
choose a seed because the test passes with it.

Randomness is not the only cause of flaky tests. Tests that depend on
timing, on the network, on the order in which tests run, or on parallel
reductions that sum in varying order fail at random as well
(Subsection~\ref{sec:isolated}).



\subsection{Most bugs live at the edges}
\label{sec:edges}

The tests of every entry point should cover the edges of its input space,
not only typical inputs, because code breaks at the edges and typical test
data never go there. Edges include:

\begin{itemize}
  \item \textbf{Sizes.} Empty input, a single element, a single group or
    category, and sizes that do not divide evenly into chunks or batches.
  \item \textbf{Values.} Zeros, ties, duplicates, columns with zero
    variance, and values at the boundary of a parameter's range, such as a
    probability of zero or one.
  \item \textbf{Missing and special values.} NA, NaN, infinities, and masked
    entries.
  \item \textbf{Magnitudes.} Very small and very large values, and nearly
    singular or ill-conditioned inputs (\DesignRef{Prefer numerically stable
    formulations}).
  \item \textbf{Structure.} Unsorted input, and identifiers that look like
    numbers or contain spaces.
\end{itemize}

At each edge, test the documented behavior: a correct result, or the
documented error or warning (\DesignRef{Raise, warn, or stay silent}). The
policy for missing values of each entry point (\DesignRef{Mind the gap,
don't fill it}) is tested this way, with missing values in every position
that the policy treats differently. Where nothing is documented for an edge,
writing the test reveals the gap in the documentation.



\subsection{Let the computer pick the inputs}
\label{sec:generated-inputs}

People choose test inputs that they expect to work. Where an invariant
(Subsection~\ref{sec:invariants}) or a metamorphic relation
(Subsection~\ref{sec:metamorphic}) holds for a whole domain of inputs, test
it on generated inputs as well as on chosen ones. Property-based testing
frameworks, such as Hypothesis for Python and quickcheck for R, generate
hundreds of inputs per test, search the edges deliberately
(Subsection~\ref{sec:edges}), and reduce a failing input to a minimal one,
which makes the bug easy to see.

Generate inputs from the whole domain that the function documents. A domain
narrowed until the tests pass excludes exactly the inputs the test was meant
to find.

Generated inputs are random, so Subsection~\ref{sec:flaky-tests} applies:
in the test suite, they must be the same on every run. Search for new
failing inputs separately, for example with the slow tests
(Subsection~\ref{sec:small-tests}), and turn every failing input the search
finds into a fixed test (Subsection~\ref{sec:bug-tests}).



\subsection{Many small tests, few big ones}
\label{sec:small-tests}

Most tests should be unit tests: each runs one function on a tiny input and
finishes in a fraction of a second. A suite of such tests runs in minutes,
so it is run before every change rather than avoided, and a failing unit
test points to the function that broke. A few end-to-end tests then run
whole entry points on a small dataset, to show that the parts fit together;
for workflow code, such a test is required (\DesignRef{Test the whole
workflow on a small dataset}).

Tests that are slow by nature, such as studies of coverage
(Subsection~\ref{sec:recovery}), convergence on fine grids
(Subsection~\ref{sec:convergence-order}), or comparisons on large inputs,
should be marked as slow and run separately, on a schedule and before every
release, rather than dropped. Slowness is no reason to skip a test silently
(Subsection~\ref{sec:skips}). Performance at the scale of real applications
is the subject of benchmarks, not of tests (\DesignRef{Size matters}).



\subsection{Every test stands on its own}
\label{sec:isolated}

Each test should set up everything it needs and leave nothing behind, so
that it gives the same outcome alone, in any order, and in parallel with
other tests. A test that relies on a seed, an option, or a file left by
another test passes or fails depending on which tests ran before it, and a
test that changes global state, such as the random generator, the working
directory, or display options, breaks others (\DesignRef{Do not touch global
state}). Write files only to a temporary directory that the test runner
provides and removes.

Tests should depend on nothing outside the repository and the declared test
dependencies, such as the network, a database, a particular machine, or the
current date. Such dependencies make tests fail for reasons unrelated to the
code (Subsection~\ref{sec:flaky-tests}), and the outcome of a test run
cannot be reproduced later.

Keep test data small and in the repository. Prefer synthetic data generated
in the test from a fixed seed, which documents how the data were made and
keeps them next to the claim they test. Where real data matter, for example
for a file format, use an excerpt small enough for the repository and
licensed for redistribution. Confidential data, such as patient records,
must not be part of a test suite: a repository is copied to every machine
and every person that works on the code, and its history keeps every file
ever committed.



\subsection{Unchanged is not correct}
\label{sec:snapshots}

A snapshot test compares a result with a reference output stored from an
earlier run of the code. It catches silent changes, such as numbers that
move after a refactoring or an update of a dependency, and it is often the
only practical test of a large output. But it shows only that the result is
unchanged, not that it is correct: a snapshot of a wrong result protects the
error.

Store a reference output only after checking it against an oracle
(Section~\ref{sec:sources}), and record how it was made: the version of the
package and the parameters (\DesignRef{Make every result traceable to a
version}). Compare it within a tolerance (Subsection~\ref{sec:float-tests}),
not bit for bit.

A failing snapshot test means that a result changed, and someone has to
decide whether the change is a fix, an intended change of numerical
behavior, or a bug. The reference output must not be regenerated just to
make the suite pass. Each update should state why the result changed, and a
change of numerical behavior goes into the changelog (\DesignRef{Design
stable APIs}).



% ==========================================================
\section{Keeping tests honest}
\label{sec:honest}
% ==========================================================

\subsection{Every bug becomes a test}
\label{sec:bug-tests}

Every fixed bug should gain a test that fails without the fix. Write the test
first: reproduce the bug, see the test fail, then fix the code and see the
test pass. A test written after the fix may pass for reasons unrelated to
the bug, and only its failure shows that it detects the bug.

A bug is also evidence about the tests. Ask which oracle would have caught
it and why none did, and look for the same mistake elsewhere in the code: a
wrong index or a wrong convention is rarely made only once.



\subsection{A test that cannot fail tests nothing}
\label{sec:can-fail}

Every new test should be seen to fail at least once. A test that has only
ever passed may never have checked anything. Common reasons are an
assertion that is never reached, for example inside a loop over an empty
list; a result compared with itself, or with an expected value from the same
code (Subsection~\ref{sec:oracle}); a tolerance so loose that any result
passes; an error caught and discarded inside the test; and a test that the
test runner never collects, for example because its file is misnamed.

Break the code on purpose and check that a test fails: flip a sign, drop the
last element, swap two axes, or change a constant by a few percent. Mutation
testing tools, available for several host stacks, make such changes
automatically and report those that no test detected.

Code coverage, the share of the code that the tests run, shows what is not
tested, not what is. Code that no test runs is untested, but code that a
test runs is tested only as far as the test checks its result.



\subsection{Skipped is not passed}
\label{sec:skips}

A test that does not run checks nothing, and a test report that hides it
looks like success (\DesignRef{Better loud than wrong}). Tests that are
skipped, for example because an optional dependency or an accelerator is
missing, and tests marked as expected failures must therefore appear as such
in the test report, with their reason.

Every test should run in at least one configuration of the automatic test
runs. A test that is skipped everywhere leaves its code path unexercised,
and such a path is not supported (\DesignRef{Testing is not optional}):
make the test run, or stop claiming support. A failing test must not be
deleted, disabled, or given a looser tolerance just to make the suite pass.
A known bug that cannot be fixed yet may be marked as an expected failure
that refers to its bug report, so that the suite shows when it is fixed.



\subsection{Tests are code too}
\label{sec:test-code}

Test code follows the same principles as the code it tests: it is read,
reviewed, and maintained, and it is read most carefully when a test fails.
Keep each test simple, with as little logic, branching, and looping as
possible; a test complex enough to need tests of its own cannot be trusted
to check anything. Define shared knowledge once, such as tolerances and the
generation of test data (\DesignRef{DRY: Don't repeat yourself}).

Each test should check one claim, and its name should state that claim, such
as that an estimate does not depend on the order of the rows, rather than a
number. When a test fails, its message should show the expected value, the
result, and the tolerance, so that the failure can be understood without
running the test again.



\end{document}
